\chapter*{Introduction générale}
\addcontentsline{toc}{chapter}{Introduction générale}
\markboth{INTRODUCTION GÉNÉRALE}{}

À l'ère du numérique, les organisations produisent et collectent des volumes
de données sans précédent. Transformer ces données brutes en informations
exploitables, puis en décisions éclairées, constitue l'enjeu central de la
\textit{Business Intelligence} (BI). Pourtant, malgré la maturité des outils
du marché (tableurs avancés, plateformes comme Power BI ou Tableau), la
chaîne décisionnelle demeure largement \textbf{manuelle} : un analyste doit
comprendre le schéma des données, choisir les indicateurs clés de performance
(KPI), concevoir des visualisations, puis rédiger un rapport de synthèse.
Chaque nouveau jeu de données relance ce cycle coûteux en temps et en
expertise.

L'émergence des \textbf{grands modèles de langage} (LLM) et des
\textbf{agents d'intelligence artificielle} ouvre la voie à une automatisation
beaucoup plus poussée de cette chaîne. Toutefois, confier l'intégralité de
l'analyse à un modèle de langage se heurte à deux écueils : les modèles
\emph{hallucinent} (ils inventent des chiffres ou des colonnes inexistantes)
et peinent à manipuler de manière fiable des jeux de données volumineux. À
l'inverse, une automatisation purement algorithmique produit des analyses
rigides et dépourvues de discernement métier.

C'est dans cette tension que s'inscrit notre projet de fin d'année. Nous
proposons \textbf{\montitre}, un écosystème décisionnel qui réconcilie ces
deux mondes : un \textbf{moteur déterministe} garantit l'exactitude des
calculs et la robustesse face à n'importe quel schéma de données, tandis
qu'une \textbf{équipe d'agents IA} apporte le jugement, la curation et la
narration. L'ensemble est orchestré par le \textbf{protocole MCP}
(\textit{Model Context Protocol}), un standard récent qui expose les
capacités du système comme un ensemble d'outils contrôlés par permissions et
mobilisables par les agents.

L'originalité de l'approche tient à son caractère \textbf{généraliste} : le
système ne fait \emph{aucune} hypothèse sur le domaine métier. Qu'il s'agisse
de ventes, de logistique, de ressources humaines ou de données sportives, la
plateforme découvre seule la structure des données, en déduit les indicateurs
et visualisations pertinents, puis produit un tableau de bord interactif et un
rapport décisionnel complet.

\section*{Problématique}
Comment concevoir une plateforme de Business Intelligence capable d'analyser
\textbf{automatiquement} et de manière \textbf{fiable} n'importe quel jeu de
données tabulaire, en combinant la rigueur d'un moteur déterministe et le
discernement d'agents fondés sur des modèles de langage, tout en garantissant
la traçabilité, la maîtrise des coûts et la véracité des résultats produits~?

\section*{Objectifs}
Le projet poursuit les objectifs suivants :
\begin{itemize}
  \item concevoir une architecture \textbf{multi-agents} orchestrée par un
        serveur MCP, avec un contrôle d'accès par rôle ;
  \item développer un \textbf{moteur adaptatif} de profilage, de calcul
        d'indicateurs et de génération de visualisations, indépendant du
        domaine métier ;
  \item produire automatiquement un \textbf{tableau de bord interactif} et un
        \textbf{rapport décisionnel} (web et PDF) reproductibles ;
  \item intégrer des mécanismes de \textbf{fiabilité} : filets de sécurité,
        vérification des chiffres, suivi des coûts ;
  \item valider l'approche par une \textbf{étude de cas} réelle.
\end{itemize}

\section*{Organisation du rapport}
Ce rapport s'articule autour de quatre chapitres. Le
\textbf{chapitre~\ref{chap:contexte}} présente le contexte général, la
problématique et l'état de l'art (BI, LLM, agents, protocole MCP). Le
\textbf{chapitre~\ref{chap:conception}} détaille l'analyse des besoins et la
conception de l'architecture. Le \textbf{chapitre~\ref{chap:realisation}}
décrit la réalisation et l'implémentation des différents modules. Enfin, le
\textbf{chapitre~\ref{chap:tests}} expose la stratégie de tests, les
résultats obtenus et une étude de cas approfondie, avant la conclusion et les
perspectives.
